\section{One affine group for the three \{6,4\} families}\label{one-affine-group-for-the-three-64-families}

Define integral affine transformations on \(\mathbb R^3\) by

\[a(x,y,z)=(y,x,z),\quad b(x,y,z)=(1-x,z,y),\quad
c(x,y,z)=(x,y,-z).\tag{5}\]

They satisfy the complete source relators

\[U=\langle a,b,c\mid a^2,b^2,c^2,(ab)^6,(bc)^4,(ac)^2,
[a,(bc)^2]^2\rangle.\tag{6}\]

Their linear parts are signed permutation matrices. On the signed basis
\((e_1,e_2,e_3,-e_1,-e_2,-e_3)\) they act as

\[ (1\ 2)(4\ 5),\quad(1\ 4)(2\ 3)(5\ 6),\quad(3\ 6).\tag{7}\]

They generate the full signed permutation group \(Q\), of order 48:
the coordinate permutations generate \(S_3\), and conjugates of the last
sign change generate all coordinate sign changes. The exact Schreier
transition calculation of Lemma AL.1 gives translation kernel

\[L=2\mathbb Z^3+\mathbb Z(1,1,1),\qquad[\mathbb Z^3:L]=4.\tag{8}\]

For example the transitions include vectors generating \(2\mathbb Z^3\)
and an all-odd vector; every transition has three coordinates of the
same parity. These two containments prove (8). A separate exact
Reidemeister--Schreier/Tietze computation for the kernel of the marked
map (6) to (7) returns three generators with the three pairwise
commutator relations. Its actual presentation is free abelian of rank
three. Lemma AL.2 therefore identifies the full presentation (6) with
the affine group (5). The computation proves a statement about the
abstract kernel, not just about its abelianization.

Here are the source's three triples of translation words. We keep them
because matching their actual values is necessary to recover the source
presentations.

\[\begin{array}{lll}
x_1=(b c^{ba})^2,&y_1=(c^b b^a)^2,&z_1=(c b^{ab})^2,\\
x_2=(a^b b^{ac})^2,&y_2=(b^{acb}a)^2,&z_2=(a^{bcb}b^{ab})^2,\\
x_3=(bac)^6,&y_3=(cba)^6,&z_3=((cba)^b)^6.
\end{array}\tag{9}\]

Their translation vectors are, in the order \(x_i,y_i,z_i\),

\[\begin{aligned}
i=1:&\quad (2,0,0),\ (0,-2,0),\ (0,0,-2),\\
i=2:&\quad (2,0,2),\ (2,-2,0),\ (-2,-2,0),\\
i=3:&\quad (2,-2,2),\ (2,-2,-2),\ (-2,-2,-2).
\end{aligned}\]

The invariant period lattices and their indices are

{
\begin{longtable}[]{@{}llr@{}}
\toprule\noalign{}
Triple & Lattice \(\Lambda_i\) & \([L:\Lambda_i]\) \\
\midrule\noalign{}
\endhead
\bottomrule\noalign{}
\endlastfoot
1 & \(2\mathbb Z^3\) & 2 \\
2 & \(2D_3\) & 4 \\
3 & \(4\mathbb Z^3+2\mathbb Z(1,1,1)\) & 8 \\
\end{longtable}
}

Here \(D_3=\{(u,v,w)\in\mathbb Z^3:u+v+w\text{ is even}\}\). Their
basis determinants have absolute values \(8,16,32\), respectively, and
all three lattices are invariant under signed permutations. Imposing
\(x_i^m=y_i^m=z_i^m=1\) gives exactly \(U/(m\Lambda_i)\), by Theorem AL.3.
Its order is \(96m^3,192m^3\) or \(384m^3\).

The adjacent affine subgroups have orders twelve and eight, and inject
into their linear images. The latter meet only in \(1\) and the linear
part of \(b\). This is checked on the finite point group and implies IP
in every period quotient. The adjacent orders remain six and four.
The point group \(C_2^3\rtimes S_3\) is solvable, so every quotient is
solvable. This proves the full source Theorem 4.3. Taking \(m=2^j\) gives
exponents \(5+3j,6+3j,7+3j\), proving its entire Corollary 4.4 without
three separate atlas identifications.

The matrix convention is explicit. The affine products in (5) act on
the rightmost factor first. GAP uses right point actions; a point
permutation corresponds to the inverse column-action matrix. Since
the three generating matrices are involutions, their displayed
permutations agree with (7), and the kernel calculation is unchanged.
